\documentclass[11pt]{article}
\usepackage[utf8]{inputenc}
\usepackage[T1]{fontenc}
\usepackage{amsmath}
\usepackage{amssymb}
\usepackage{multicol}
\usepackage{geometry}
\usepackage{tikz}
\usepackage{enumitem}
\usepackage{xcolor}
\usepackage{titlesec}

% Configurações de layout
\geometry{a4paper, left=1cm, right=1cm, top=0.5cm, bottom=1.2cm}
\setlength{\columnseprule}{0.4pt}

% Cores para títulos
\titleformat{\section}{\normalfont\Large\bfseries\color{blue}}{\thesection}{1em}{}
\titleformat{\subsection}{\normalfont\large\bfseries\color{red}}{\thesubsection}{1em}{}

% Título principal
\title{\textcolor{red}{Gabarito Comentado: Função do Segundo Grau
    \\Primeiro Trimestre}}
\author{Jefferson}
\date{}

\begin{document}

\maketitle
\vspace{-1cm}

\begin{multicols}{2}

% ============================================================
% GABARITO: EXERCÍCIOS BÁSICOS - RAÍZES
% ============================================================
\section*{Gabarito: Raízes}

\subsection*{Questão 1}
\textbf{Enunciado:} Calcule as raízes de $f(x) = x^2 - 7x + 12$.

\textbf{Solução:} Usando a fórmula de Bhaskara com $a = 1$, $b = -7$, $c = 12$:
\[
\Delta = (-7)^2 - 4(1)(12) = 49 - 48 = 1
\]
\[
x = \frac{7 \pm \sqrt{1}}{2} = \frac{7 \pm 1}{2}
\]
Logo, $x_1 = 4$ e $x_2 = 3$.

\subsection*{Questão 2}
\textbf{Enunciado:} Determine o discriminante de $f(x) = 2x^2 - 3x + 1$ e classifique as raízes.

\textbf{Solução:} Com $a = 2$, $b = -3$, $c = 1$:
\[
\Delta = (-3)^2 - 4(2)(1) = 9 - 8 = 1 > 0
\]
Como $\Delta > 0$, a função possui \textbf{duas raízes reais e distintas}.

\subsection*{Questão 3}
\textbf{Enunciado:} Encontre as raízes de $f(x) = x^2 - 2x - 8$.

\textbf{Solução:} Com $a = 1$, $b = -2$, $c = -8$:
\[
\Delta = (-2)^2 - 4(1)(-8) = 4 + 32 = 36
\]
\[
x = \frac{2 \pm \sqrt{36}}{2} = \frac{2 \pm 6}{2}
\]
Logo, $x_1 = 4$ e $x_2 = -2$.

\subsection*{Questão 4}
\textbf{Enunciado:} Qual é o valor de $\Delta$ para $f(x) = x^2 + x + 1$? A função possui raízes reais?

\textbf{Solução:} Com $a = 1$, $b = 1$, $c = 1$:
\[
\Delta = 1^2 - 4(1)(1) = 1 - 4 = -3 < 0
\]
Como $\Delta < 0$, a função \textbf{não possui raízes reais}.

\subsection*{Questão 5}
\textbf{Enunciado:} Resolva $3x^2 - 12x + 9 = 0$.

\textbf{Solução:} Simplificando por 3: $x^2 - 4x + 3 = 0$. Com $a = 1$, $b = -4$, $c = 3$:
\[
\Delta = 16 - 12 = 4
\]
\[
x = \frac{4 \pm 2}{2} \Rightarrow x_1 = 3, \, x_2 = 1
\]

% ============================================================
% GABARITO: EXERCÍCIOS - VÉRTICE
% ============================================================
\section*{Gabarito: Vértice}

\subsection*{Questão 6}
\textbf{Enunciado:} Calcule as coordenadas do vértice de $f(x) = x^2 - 4x + 3$.

\textbf{Solução:} Com $a = 1$, $b = -4$:
\[
x_v = -\frac{-4}{2(1)} = 2, \quad y_v = f(2) = 4 - 8 + 3 = -1
\]
Vértice: $V = (2, -1)$.

\subsection*{Questão 7}
\textbf{Enunciado:} Determine o vértice e diga se é máximo ou mínimo: $f(x) = -x^2 + 6x - 5$.

\textbf{Solução:} Com $a = -1$ (concavidade para baixo), $b = 6$:
\[
x_v = -\frac{6}{2(-1)} = 3, \quad y_v = f(3) = -9 + 18 - 5 = 4
\]
\begin{enumerate}[label=\alph*)]
\item Vértice: $(3, 4)$
\item \textbf{Máximo} (pois $a < 0$)
\item Eixo de simetria: $x = 3$
\end{enumerate}

\subsection*{Questão 8}
\textbf{Enunciado:} Encontre o vértice de $f(x) = 2x^2 - 8x + 6$.

\textbf{Solução:} Com $a = 2$, $b = -8$:
\[
x_v = -\frac{-8}{2(2)} = 2, \quad y_v = f(2) = 8 - 16 + 6 = -2
\]
Vértice: $V = (2, -2)$.

\subsection*{Questão 9}
\textbf{Enunciado:} Qual é o valor máximo da função $f(x) = -3x^2 + 12x - 9$?

\textbf{Solução:} Com $a = -3$, $b = 12$:
\[
x_v = -\frac{12}{2(-3)} = 2, \quad y_v = f(2) = -12 + 24 - 9 = 3
\]
Valor máximo: $\mathbf{3}$.

\subsection*{Questão 10}
\textbf{Enunciado:} Para $f(x) = x^2 - 5x + 4$, determine a abscissa do vértice.

\textbf{Solução:} Com $a = 1$, $b = -5$:
\[
x_v = -\frac{-5}{2(1)} = 2,5
\]

% ============================================================
% GABARITO: EXERCÍCIOS - FORMAS
% ============================================================
\section*{Gabarito: Formas Alternativas}

\subsection*{Questão 11}
\textbf{Enunciado:} Converta $f(x) = (x - 3)(x + 2)$ para a forma geral.

\textbf{Solução:} Expandindo:
\[
f(x) = x^2 + 2x - 3x - 6 = x^2 - x - 6
\]

\subsection*{Questão 12}
\textbf{Enunciado:} Escreva $f(x) = x^2 - 6x + 8$ na forma fatorada.

\textbf{Solução:} Encontrando as raízes: $\Delta = 36 - 32 = 4$, então $x = \frac{6 \pm 2}{2} \Rightarrow x_1 = 4, x_2 = 2$.
\begin{enumerate}[label=\alph*)]
\item Raízes: $x = 2$ e $x = 4$
\item Forma fatorada: $f(x) = (x - 2)(x - 4)$
\end{enumerate}

\subsection*{Questão 13}
\textbf{Enunciado:} Transforme $f(x) = x^2 - 4x + 4$ na forma canônica.

\textbf{Solução:} Observe que $x^2 - 4x + 4 = (x - 2)^2$. Logo:
\[
f(x) = (x - 2)^2 + 0 = (x - 2)^2
\]
Forma canônica: $f(x) = 1 \cdot (x - 2)^2 + 0$.

\subsection*{Questão 14}
\textbf{Enunciado:} Dada $f(x) = 2(x - 1)^2 - 3$ (forma canônica), identifique o vértice.

\textbf{Solução:} Na forma $f(x) = a(x - x_v)^2 + y_v$, temos $x_v = 1$ e $y_v = -3$.
Vértice: $(1, -3)$.

\subsection*{Questão 15}
\textbf{Enunciado:} Expanda $(x - 2)^2 + 5$ para obter a forma geral.

\textbf{Solução:} 
\[
f(x) = x^2 - 4x + 4 + 5 = x^2 - 4x + 9
\]

% ============================================================
% GABARITO: EXERCÍCIOS - SINAL
% ============================================================
\section*{Gabarito: Sinal e Inequações}

\subsection*{Questão 16}
\textbf{Enunciado:} Estude o sinal de $f(x) = x^2 - 5x + 6$.

\textbf{Solução:} Raízes: $x_1 = 2, x_2 = 3$. Como $a = 1 > 0$:
\begin{itemize}
\item $f(x) > 0$ se $x < 2$ ou $x > 3$
\item $f(x) = 0$ se $x = 2$ ou $x = 3$
\item $f(x) < 0$ se $2 < x < 3$
\end{itemize}

\subsection*{Questão 17}
\textbf{Enunciado:} Para quais valores de $x$ tem-se $x^2 - 4 > 0$?

\textbf{Solução:}
\begin{enumerate}[label=\alph*)]
\item Raízes: $x^2 - 4 = 0 \Rightarrow x = \pm 2$
\item Sinal: Como $a = 1 > 0$, a parábola é positiva fora das raízes
\item Solução: $x < -2$ ou $x > 2$, ou em notação de intervalo: $(-\infty, -2) \cup (2, \infty)$
\end{enumerate}

\subsection*{Questão 18}
\textbf{Enunciado:} Resolva a inequação $-x^2 + 4x - 3 \geq 0$.

\textbf{Solução:} Raízes de $-x^2 + 4x - 3 = 0$: $x = 1$ e $x = 3$. Como $a = -1 < 0$:
\[
-x^2 + 4x - 3 \geq 0 \Rightarrow 1 \leq x \leq 3
\]

\subsection*{Questão 19}
\textbf{Enunciado:} Estude o sinal de $f(x) = 2x^2 - 8x + 8$.

\textbf{Solução:} Simplificando: $f(x) = 2(x^2 - 4x + 4) = 2(x - 2)^2$. Como $\Delta = 0$:
\begin{itemize}
\item $f(x) > 0$ para todo $x \neq 2$
\item $f(x) = 0$ se $x = 2$
\item $f(x) < 0$ nunca
\end{itemize}

\subsection*{Questão 20}
\textbf{Enunciado:} Para que valores de $x$ temos $f(x) = x^2 - 2x - 3 < 0$?

\textbf{Solução:} Raízes: $x = 3$ e $x = -1$. Como $a = 1 > 0$:
\[
x^2 - 2x - 3 < 0 \Rightarrow -1 < x < 3
\]

% ============================================================
% GABARITO: EXERCÍCIOS - APLICAÇÕES
% ============================================================
\section*{Gabarito: Exercícios de Aplicação}

\subsection*{Questão 21}
\textbf{Enunciado:} Um projétil é lançado e sua altura é $h(t) = -5t^2 + 20t$. Qual é a altura máxima?

\textbf{Resposta:} $t_v = -\frac{20}{2(-5)} = 2$ segundos. $h_{max} = h(2) = -5(4) + 20(2) = 20$ metros.

\subsection*{Questão 22}
\textbf{Enunciado:} Lucro: $L(x) = -2x^2 + 100x - 1200$. Quantidade que maximiza?

\textbf{Resposta:} $x_v = -\frac{100}{2(-2)} = 25$ unidades.

\subsection*{Questão 23}
\textbf{Enunciado:} Terreno retangular com $A(x) = x(20-x)$. Área máxima?

\textbf{Resposta:} $A(x) = -x^2 + 20x$. $x_v = 10$ metros. $A_{max} = 100$ m².

\subsection*{Questão 24}
\textbf{Enunciado:} Temperatura: $T(h) = -h^2 + 24h - 80$. Máxima?

\textbf{Resposta:} $h_v = 12$ horas. $T_{max} = -144 + 288 - 80 = 64$ graus.

\subsection*{Questão 25}
\textbf{Enunciado:} Custo: $C(x) = x^2 - 10x + 50$. Mínimo?

\textbf{Resposta:} $x_v = 5$ unidades. $C_{min} = 25 - 50 + 50 = 25$.

\subsection*{Questão 26}
\textbf{Enunciado:} Preço: $P(x) = -2x^2 + 100x + 500$. Máximo em qual mês?

\textbf{Resposta:} $x_v = -\frac{100}{2(-2)} = 25$ meses.

\subsection*{Questão 27}
\textbf{Enunciado:} População: $P(t) = -t^2 + 12t + 50$. Máxima?

\textbf{Resposta:} $t_v = 6$ unidades de tempo. $P_{max} = -36 + 72 + 50 = 86$.

\subsection*{Questão 28}
\textbf{Enunciado:} Receita: $R(x) = -3x^2 + 150x$. Quantidade que maximiza?

\textbf{Resposta:} $x_v = -\frac{150}{2(-3)} = 25$ unidades.

\subsection*{Questão 29}
\textbf{Enunciado:} Altura da bola: $h(t) = -10t^2 + 30t$. Quando retorna ao solo?

\textbf{Solução:} $h(t) = 0 \Rightarrow -10t^2 + 30t = 0 \Rightarrow t(-10t + 30) = 0$.
Logo, $t = 0$ (lançamento) ou $t = 3$ segundos (retorno).
\textbf{Resposta:} 3 segundos.

\subsection*{Questão 30}
\textbf{Enunciado:} Eficiência: $E(x) = -x^2 + 8x - 12$. Máxima?

\textbf{Solução:} $x_v = -\frac{8}{2(-1)} = 4$. $E_{max} = -16 + 32 - 12 = 4$.
\textbf{Resposta:} Eficiência máxima é 4.

\end{multicols}

\end{document}
